\part{The QFT and phase estimation}\label{part:qft}

\begin{frame}
  \partpage
\end{frame}

%%%%%%%%%%%%%%%%%%%%%%%%%%%%%%%%%%%%%%%%%%%%%%%%%%%%%%%%%%%%%%%%%%%%%%%%%%%%

%% Quantum Phase Estimation

%%%%%%%%%%%%%%%%%%%%%%%%%%%%%%%%%%%%%%%%%%%%%%%%%%%%%%%%%%%%%%%%%%%%%%%%%%%%

\begin{frame}{Quantum phase estimation}

\begin{problem}
We are given a unitary $U$ and an eigenvector $|\psi\>$ of $U$ with unknown eigenvalue \pause

\bigskip		
We seek to determine its eigenphase $\varphi\in[0,1)$ such that 
\[
   U|\psi\> = e^{2 \pi i \varphi}|\psi\>
\] \pause

\medskip
More precisely, we want to obtain an estimate $\hat{\varphi}$ such that
\[
   \Pr\big(|\hat{\varphi} - \varphi| \le \frac{1}{2^n}\big) \ge \frac{3}{4}
\]
The deviation $|\hat{\varphi} - \varphi|$ is computed modulo $1$ 
\end{problem}

\end{frame}

%%%%%%%%%%%%%%%%%%%%%%%%%%%%%%%%%%%%%%%%%%%%%%%%%%%%%%%%%%%%%%%

%% Phase Kick Back

%%%%%%%%%%%%%%%%%%%%%%%%%%%%%%%%%%%%%%%%%%%%%%%%%%%%%%%%%%%%%%%

\begin{frame}{Phase kick back}

\hspace{4cm}
\Qcircuit @C=1em @R=0.7em {
\lstick{|+\>}    & \ctrl{1} & \qw  \\
\lstick{|\psi\>} & \gate{U} & \qw 
}

\bigskip
\begin{eqnarray*}
\frac{|0\> + |1\>}{\sqrt{2}} \otimes |\psi\>  \pause
&  = &
\frac{|0\>}{\sqrt{2}} \otimes |\psi\> + \frac{|1\>}{\sqrt{2}} \otimes |\psi\> \\ \pause
& \mapsto &
\frac{|0\>}{\sqrt{2}} \otimes |\psi\> + \frac{|1\>}{\sqrt{2}} \otimes U |\psi\> \\ \pause
& = &
\frac{|0\>}{\sqrt{2}} \otimes |\psi\> + \frac{|1\>}{\sqrt{2}} \otimes e^{2\pi i\varphi} |\psi\> \\ \pause
& = &
\frac{|0\>}{\sqrt{2}} \otimes |\psi\> + \frac{e^{2\pi i\varphi} |1\>}{\sqrt{2}} \otimes |\psi\> \\ \pause
& = &
\frac{|0\> + e^{2\pi i\varphi} |1\>}{\sqrt{2}} \otimes |\psi\> 
\end{eqnarray*}

\end{frame}

%%%%%%%%%%%%%%%%%%%%%%%%%%%%%%%%%%%%%%%%%%%%%%%%%%%%%%%%%%%%%%%

%% Phase Kick Back

%%%%%%%%%%%%%%%%%%%%%%%%%%%%%%%%%%%%%%%%%%%%%%%%%%%%%%%%%%%%%%%

\begin{frame}{Phase kick back}

\hspace{4cm}
\Qcircuit @C=1em @R=0.7em {
\lstick{|+\>}    & \ctrl{1} & \qw  \\
\lstick{|\psi\>} & \gate{U} & \qw 
}

\bigskip
\[
\frac{|0\> + |1\>}{\sqrt{2}} \otimes |\psi\>  \pause
\mapsto
\frac{|0\> + e^{2\pi i\varphi} |1\>}{\sqrt{2}} \otimes |\psi\> 
\] \pause

\medskip
The eigenstate $|\psi\>$ in the {\color{blue} target} register emerges {\color{blue} unchanged} \pause

\bigskip
$\Rightarrow$ It suffices to focus on the control register \pause

\bigskip
The state $|0\>+|1\>$ of the {\color{red} control} qubit is {\color{red} changed} to $|0\>+e^{2\pi i \varphi}|1\>$ by phase kick back 

\end{frame}

%%%%%%%%%%%%%%%%%%%%%%%%%%%%%%%%%%%%%%%%%%%%%%%%%%%%%%%%%%%%%%%

%% Phase Kick Back

%%%%%%%%%%%%%%%%%%%%%%%%%%%%%%%%%%%%%%%%%%%%%%%%%%%%%%%%%%%%%%%

\begin{frame}{Hadamard test + phase kick back}

\hspace{4cm}
\Qcircuit @C=1em @R=0.7em {
\lstick{|0\>}    & \gate{H}    & \ctrl{1} & \gate{H} & \meter \\
\lstick{|\psi\>} & {/}\qw      & \gate{U} & \qw      & \qw
}

\bigskip
\begin{eqnarray*}
&& 
\frac{|0\> + e^{2\pi i\varphi} |1\>}{\sqrt{2}} \\ \pause
& \mapsto &
\frac{1}{2}\, \big( (|0\> + |1\>) + e^{2\pi i\varphi} ( |0\> - |1\> ) \big) \\ \pause
& \mapsto &
\frac{1}{2}\, \big( ( 1 + e^{2\pi i\varphi}) |0\> + ( 1 - e^{2\pi i\varphi}) |1\> ) \big)  := |\varphi\> 
\end{eqnarray*}

\end{frame}

%%%%%%%%%%%%%%%%%%%%%%%%%%%%%%%%%%%%%%%%%%%%%%%%%%%%%%%%%%%%%%%

%% Phase Kick Back

%%%%%%%%%%%%%%%%%%%%%%%%%%%%%%%%%%%%%%%%%%%%%%%%%%%%%%%%%%%%%%%

\begin{frame}{Hadamard test + phase kick back}

\[
|\varphi\> = \frac{1}{2}\, \big( ( 1 + e^{2\pi i\varphi}) |0\> + ( 1 - e^{2\pi i\varphi}) |1\> ) \big) \pause
\]

\bigskip
The probability of obtaining $0$ is \pause

\medskip
\begin{eqnarray*}
\Pr(0) 
& = & 
\| |0\>\<0| \, |\varphi\> \|^2 \\ \pause
& = &
| \frac{1}{2} \, (1 + e^{2\pi i \varphi}) |^2 \\ \pause
& = &
\frac{1}{4} \, |e^{\pi i\varphi} + e^{-\pi i\varphi}) |^2 \\ \pause
& = &
\frac{1}{4} \, | 2 \cos(\pi\varphi)|^2 \\ \pause
& = &
\cos^2(\pi\varphi) = \frac{1}{2} \, \big(1 +\cos(2\pi\varphi)\big)
\end{eqnarray*}
\end{frame}

%%%%%%%%%%%%%%%%%%%%%%%%%%%%%%%%%%%%%%%%%%%%%%%%%%%%%%%%%%%%%%%%%%%%%%%%%%%%%%%%%%%%%%%%%%%%%%%%%%%%%%%%%%%%

%% Phase kick back 

%%%%%%%%%%%%%%%%%%%%%%%%%%%%%%%%%%%%%%%%%%%%%%%%%%%%%%%%%%%%%%%%%%%%%%%%%%%%%%%%%%%%%%%%%%%%%%%%%%%%%%%%%%%%
\begin{frame}{Phase kick back due to higher powers of $U$}

	For arbitrary $k$, we obtain
	
	\bigskip
	
	\bigskip
	\hspace{1cm}
	\Qcircuit @C=1.0em @R=1.0em {  
		\lstick{|0\>}    & \qw  &\gate{H}  &  \ctrl{1}       &  \qw  &  \qw &  \qw &  \rstick{\frac{1}{\sqrt{2}}({|0\> + e^{2\pi i 2^k\varphi}|1\>})}\\
	 	\lstick{|\psi\>} & \qw  &{/}\qw    &  \gate{U^{2^k}} &  \qw  &  \qw &  \qw &  \rstick{|\psi\>}
	}			 	
	
	\bigskip
	
	since
	\begin{eqnarray*}
		U^{2^k}|\psi\>=e^{2\pi i 2^k \varphi}|\psi\>
	\end{eqnarray*}

\end{frame}

%%%%%%%%%%%%%%%%%%%%%%%%%%%%%%%%%%%%%%%%%%%%%%%%%%%%%%%%%%%%%%%%%%%%%%%%%%%%

%% Phase estimation

%%%%%%%%%%%%%%%%%%%%%%%%%%%%%%%%%%%%%%%%%%%%%%%%%%%%%%%%%%%%%%%%%%%%%%%%%%%%

\begin{frame}{Phase kick back part of phase estimation}

\hspace{1cm}
\Qcircuit @C=1.0em @R=0.7em {
\lstick{|0\>}    & \qw & \gate{H} & \qw      & \qw      & \qw &        & & \qw & \ctrl{4} & \qw & \qw & \rstick{\frac{|0\> + e^{2\pi i 2^{n-1}\varphi}|1\>}{\sqrt{2}}} \\
\lstick{|0\>}    & \qw & \gate{H} & \qw      & \ctrl{3} & \qw &        & & \qw & \qw      & \qw & \qw & \rstick{\frac{|0\> + e^{2\pi i 2^{n-2}\varphi}|1\>}{\sqrt{2}}} \\
                 &     & \raisebox{0.5em}{\vdots}   &          &          &     & \cdots & &     &          &     &     & \rstick{\raisebox{0.5em}{\vdots}} \\
\lstick{|0\>}    & \qw & \gate{H} & \ctrl{1} & \qw      & \qw &        & & \qw & \qw      & \qw & \qw & \rstick{\frac{|0\> + e^{2 \pi i 2^0\varphi}|1\>}{\sqrt{2}}} \\
%
\lstick{|\psi\>} & \qw & {/} \qw  & \gate{U^{2^0}} & \gate{U^{2^1}} & \qw & & & \qw & \gate{U^{2^{n-1}}} & \qw & \qw & \rstick{|\psi\>}
}

\bigskip

\bigskip
We set
\[
|\varphi\> := 
\frac{|0\> + e^{2 \pi i 2^{n-1} \varphi}|1\>}{\sqrt{2}} \!\otimes\!
\frac{|0\> + e^{2 \pi i 2^{n-2} \varphi}|1\>}{\sqrt{2}} \!\otimes\!\cdots\!\otimes
\frac{|0\> + e^{2 \pi i 2^0 \varphi}    |1\>}{\sqrt{2}} 
\]

\end{frame}

%%%%%%%%%%%%%%%%%%%%%%%%%%%%%%%%%%%%%%%%%%%%%%%%%%%%%%%%%%%%%%%%%%%%%%%%%%%

%% Binary fractions

%%%%%%%%%%%%%%%%%%%%%%%%%%%%%%%%%%%%%%%%%%%%%%%%%%%%%%%%%%%%%%%%%%%%%%%%%%%

\begin{frame}{Binary fractions}

Assume that the eigenphase $\varphi$ is an exact $n$-bit binary fraction, i.e.,
\[
\varphi=0.x_1 x_2 \ldots x_n = \sum_{i=1}^n \frac{x_i}{2^i}
\]

For arbitrary $k\in\{0,\ldots,n-1\}$, we have
\begin{eqnarray*}
2^k\, \varphi 
& = &
x_1 x_2\ldots x_k . x_{k+1} \ldots x_n \\
& & \\
e^{2\pi i 2^k \varphi} 
&=&
e^{2\pi i (x_1x_2\ldots x_k . x_{k+1}\ldots x_n)}\\
&=&
e^{2\pi i (x_1x_2\ldots x_k + 0. x_{k+1}\ldots x_n)} \\
&=&
e^{2\pi i (x_1x_2\ldots x_k)} \cdot e^{2\pi i (0. x_{k+1}\ldots x_n)} \\
&=& 
e^{2\pi i (0. x_{k+1}\ldots x_n)}
\end{eqnarray*}

\end{frame}

%%%%%%%%%%%%%%%%%%%%%%%%%%%%%%%%%%%%%%%%%%%%%%%%%%%%%%%%%%%%%%%%%%%%%%%%%%%%

%% Phase estimation

%%%%%%%%%%%%%%%%%%%%%%%%%%%%%%%%%%%%%%%%%%%%%%%%%%%%%%%%%%%%%%%%%%%%%%%%%%%%

\begin{frame}{Phase kick back part of phase estimation}

\Qcircuit @C=1.0em @R=0.7em {
\lstick{|0\>}&\qw& \gate{H} & \qw      & \qw      & \qw &        & & \qw & \ctrl{4} & \qw & \qw & \rstick{\frac{|0\> + e^{2\pi i 0.x_n}|1\>}{\sqrt{2}}}\\
\lstick{|0\>}&\qw& \gate{H} & \qw      & \ctrl{3} & \qw &        & & \qw & \qw      & \qw & \qw & \rstick{\frac{|0\> + e^{2\pi i 0.x_{n-1} x_n}|1\>}{\sqrt{2}}}\\
             &   & \raisebox{0.5em}{\vdots}   &          &          &     & \cdots & &     &          &     &     & \rstick{\raisebox{0.5em}{\vdots}} \\
\lstick{|0\>}&\qw& \gate{H} & \ctrl{1} & \qw      & \qw &        & & \qw & \qw      & \qw & \qw & \rstick{\frac{|0\> + e^{2 \pi i 0.x_1 \ldots x_{n-1} x_n}|1\>}{\sqrt{2}}}\\
%
\lstick{|\psi\>} & \qw & {/} \qw  & \gate{U^{2^0}} & \gate{U^{2^1}} & \qw & & & \qw & \gate{U^{2^{n-1}}} & \qw & \qw & \rstick{|\psi\>}
}

\end{frame}

%%%%%%%%%%%%%%%%%%%%%%%%%%%%%%%%%%%%%%%%%%%%%%%%%%%%%%%%%%%%%%%%%%%%%%%%%%%%

%% Quantum Fourier transform

%%%%%%%%%%%%%%%%%%%%%%%%%%%%%%%%%%%%%%%%%%%%%%%%%%%%%%%%%%%%%%%%%%%%%%%%%%%%

\begin{frame}{Quantum Fourier transform}

The quantum Fourier transform $F$ is defined by

\begin{eqnarray*}
& & 
F \big( |x_n\> \otimes |x_{n-1}\> \otimes \cdots \otimes |x_1\> \big) \\ & & \\
& = &
\frac{|0\> + e^{2 \pi i 0.x_n}|1\>}{\sqrt{2}} \!\otimes\!
\frac{|0\> + e^{2 \pi i 0.x_{n-1} x_n}|1\>}{\sqrt{2}} \!\otimes\!\cdots\!\otimes
\frac{|0\> + e^{2 \pi i 0.x_1x_2\ldots x_n}|1\>}{\sqrt{2}} 
%\\ \pause
%& & \\
%& = &
%\frac{1}{\sqrt{2^n}} \sum_{y=0}^{2^n-1} e^{2\pi i \varphi y} |y\>
\end{eqnarray*}

%\medskip
%where $\varphi=0.x_1 x_2 \ldots x_n$, $y=\sum_{k=1}^n y_k 2^{n-k}$, and $|y\>=|y_n\>\otimes |y_{n-1}\>\otimes\cdots\otimes |y_1\>$ \pause

\bigskip
We use inverse quantum Fourier transform $F^\dagger$ to obtain the bits of the eigenphase \pause

\bigskip
Note: QFT is defined by $F|x_1\>\otimes|x_2\>\otimes\cdots\otimes|x_n\>=|0.x_1 x_2 \ldots x_n\>$ in the literature; we use the above definition for the sake of notational simplicity (otherwise, we would have to include the so-called bit-reversal)

\end{frame}

%%%%%%%%%%%%%%%%%%%%%%%%%%%%%%%%%%%%%%%%%%%%%%%%%%%%%%%%%%%%%%%%%%%%%%%%%%%%

%% Quantum circuit for phase estimation

%%%%%%%%%%%%%%%%%%%%%%%%%%%%%%%%%%%%%%%%%%%%%%%%%%%%%%%%%%%%%%%%%%%%%%%%%%%%

\begin{frame}{Quantum circuit for phase estimation}

\Qcircuit @C=1.0em @R=0.7em {
\lstick{|0\>}& \qw & \gate{H}                 & \qw      & \qw      & \qw &        &        & \qw & \ctrl{4} & \qw & \multigate{3}{F^\dagger} & \qw & \rstick{|x_n\>}\\
\lstick{|0\>}& \qw & \gate{H}                 & \qw      & \ctrl{3} & \qw &        &        & \qw & \qw      & \qw & \ghost{F^\dagger}        & \qw & \rstick{|x_{n-1}\>}\\
             &     & \raisebox{0.5em}{\vdots} &          &          &     & \cdots &        &     &          &     &         &     & \\
\lstick{|0\>}& \qw & \gate{H}                 & \ctrl{1} & \qw      & \qw &        &        & \qw & \qw      & \qw & \ghost{F^\dagger}        & \qw & \rstick{|x_1\>}\\
%
\lstick{|\psi\>} & \qw & {/} \qw  & \gate{U^{2^0}} & \gate{U^{2^1}} & \qw & & & \qw & \gate{U^{2^{n-1}}} & \qw & \qw & \qw & \rstick{|\psi\>}
}

\end{frame}

%%%%%%%%%%%%%%%%%%%%%%%%%%%%%%%%%%%%%%%%%%%%%%%%%%%%%%%%%%%%%%%%%%%%%%%%%%%

\begin{frame}{Inverse quantum Fourier transform for $3$ bits}

\hspace{2.0cm}
\Qcircuit @C=.4em @R=1.0em {  
\lstick{\frac{|0\> + e^{2\pi i 0.x_3}|1\>}{\sqrt{2}}}       &\qw &\gate{H} & \qw & \qw & \qw & \qw & \ctrl{1}        & \qw & \qw      & \qw & \qw & \qw & \ctrl{2}        & \qw             & \qw & \qw & \qw & \qw & \rstick{|x_3\>} \\ 
\lstick{\frac{|0\> + e^{2\pi i 0.x_2x_3}|1\>}{\sqrt{2}}}    &\qw &\qw      & \qw & \qw & \qw & \qw & \gate{R_2^\dagger} & \qw & \gate{H} & \qw & \qw & \qw & \qw             & \ctrl{1}        & \qw &\qw      &\qw & \qw & \rstick{|x_2\>} \\ 
\lstick{\frac{|0\> + e^{2\pi i 0.x_1x_2x_3}|1\>}{\sqrt{2}}} &\qw &\qw      & \qw & \qw & \qw & \qw & \qw             & \qw & \qw      & \qw & \qw & \qw & \gate{R_3^\dagger} & \gate{R_2^\dagger} & \qw &\gate{H} &\qw & \qw & \rstick{|x_1\>}  
	        \gategroup{1}{3}{1}{3}{1.3em}{--}
	        \gategroup{1}{8}{2}{10}{1.3em}{--}
	        \gategroup{1}{14}{3}{17}{1.3em}{--}
}	

\bigskip
The phase shift $R_k$ is defined by
\[
R_k :=
\left[ 
\begin{array}{cc}
 1 & 0  \\
 0 & e^{2\pi i/2^k}  \\
 \end{array} 
\right]
\]	 
\end{frame}

%%%%%%%%%%%%%%%%%%%%%%%%%%%%%%%%%%%%%%%%%%%%%%%%%%%%%%%%%%%%%%%%%%%%%%%%%%%%

%% least significant bit -- top qubit

%%%%%%%%%%%%%%%%%%%%%%%%%%%%%%%%%%%%%%%%%%%%%%%%%%%%%%%%%%%%%%%%%%%%%%%%%%%%

\begin{frame}{QPE: least significant bit -- top qubit}

\hspace{3cm}
\Qcircuit @C=1.0em @R=1.0em {  
	\lstick{\frac{|0\>+e^{2\pi i 0.x_3}|1\>}{\sqrt{2}}} &\qw &\gate{H} & \qw & \qw
	\gategroup{1}{3}{1}{3}{1.2em}{--}
}			 	

\bigskip
\begin{eqnarray*}
	&                 & \frac{|0\> + e^{2\pi i 0.x_3}|1\>}{\sqrt{2}} \pause = \frac{|0\> + (-1)^{x_3} |1\>}{\sqrt{2}} \\ \pause
	& \xrightarrow{H} & |x_3\> 
\end{eqnarray*}

\end{frame}

%%%%%%%%%%%%%%%%%%%%%%%%%%%%%%%%%%%%%%%%%%%%%%%%%%%%%%%%%%%%%%%%%%%%%%%%%%%%%

%% second bit -- middle qubit

%%%%%%%%%%%%%%%%%%%%%%%%%%%%%%%%%%%%%%%%%%%%%%%%%%%%%%%%%%%%%%%%%%%%%%%%%%%%%

\begin{frame}{QPE: second bit -- middle qubit}

\hspace{3cm}
\Qcircuit @C=1.0em @R=1.0em {
	\lstick{|x_3\>}                                           & \qw & \ctrl{1}            & \qw      & \qw & \qw \\
	\lstick{\frac{|0\> + e^{2\pi i 0.x_2 x_3}|1\>}{\sqrt{2}}} & \qw & \gate{R_2^\dagger}  & \gate{H} & \qw & \qw 
	\gategroup{1}{3}{2}{4}{1.5em}{--}
}
			              
\bigskip        
\begin{eqnarray*}
		&                                  & |x_3\> \otimes \frac{|0\> + e^{2\pi i 0.x_2 x_3}|1\>}{\sqrt{2}} \\ \pause
		& \xrightarrow{ctrl\, R_2^\dagger} & |x_3\> \otimes \frac{|0\> + e^{2\pi i 0.x_2 0}  |1\>}{\sqrt{2}} \\ \pause
    & \xrightarrow{I \otimes H}				 & |x_3\> \otimes |x_2\> 
\end{eqnarray*} 

\end{frame}

%%%%%%%%%%%%%%%%%%%%%%%%%%%%%%%%%%%%%%%%%%%%%%%%%%%%%%%%%%%%%%%%%%%%%%%%%%%%%

%% most significant bit -- bottom qubit

%%%%%%%%%%%%%%%%%%%%%%%%%%%%%%%%%%%%%%%%%%%%%%%%%%%%%%%%%%%%%%%%%%%%%%%%%%%%%

\begin{frame}{QPE: most significant bit -- bottom qubit}

\hspace{3cm}
\Qcircuit @C=1.0em @R=1.0em {
	\lstick{|x_3\>}                                        & \qw & \ctrl{2}           & \qw                &\qw      & \qw & \qw \\
	\lstick{|x_2\>}                                        & \qw & \qw                & \ctrl{1}           &\qw      & \qw & \qw \\
	\lstick{\frac{|0\> + e^{0.x_1 x_2 x_3}|1\>}{\sqrt{2}}} & \qw & \gate{R_3^\dagger} & \gate{R_2^\dagger} &\gate{H} & \qw & \qw
	\gategroup{1}{3}{3}{5}{1.5em}{--}
}

\bigskip
\begin{eqnarray*}
			& & 																	  |x_3\> \otimes |x_2\> \otimes \frac{|0\> + e^{2\pi i 0.x_1 x_2 x_3}|1\>}{\sqrt{2}}  \\ \pause
			& \xrightarrow{ctrl\, R_3^\dagger}    & |x_3\> \otimes |x_2\> \otimes \frac{|0\> + e^{2\pi i 0.x_1 x_2 0  }|1\>}{\sqrt{2}}  \\ \pause
  		& \xrightarrow{ctrl\, R_2^\dagger}    & |x_3\> \otimes |x_2\> \otimes \frac{|0\> + e^{2\pi i 0.x_1 0   0  }|1\>}{\sqrt{2}}  \\ \pause
  		& \xrightarrow{I \otimes I \otimes H} & |x_3\> \otimes |x_2\> \otimes |x_1\> 
\end{eqnarray*} 
\end{frame}

%%%%%%%%%%%%%%%%%%%%%%%%%%%%%%%%%%%%%%%%%%%%%%%%%%%%%%%%%%%%%%%%%%%%%%%%%%%%

%% Summary of phase estimation circuit

%%%%%%%%%%%%%%%%%%%%%%%%%%%%%%%%%%%%%%%%%%%%%%%%%%%%%%%%%%%%%%%%%%%%%%%%%%%%

\begin{frame}{Summary of phase estimation circuit}


  We use phase kick back due to the controlled $U^{2^k}$ gate to prepare the state
  		\[
  				\frac{|0\> + e^{2 \pi i 0.x_{k+1} x_{k+2}\ldots x_n} |1\>}{\sqrt{2}}
  		\] \pause
  
  Using the previously determined bits $x_{k+2},\ldots, x_n$, we change this state to
  	  \[
  	  		\frac{|0\> + e^{2 \pi i 0.x_{k+1} 0       \ldots 0  } |1\>}{\sqrt{2}} = \frac{|0\> + (-1)^{x_k}|1\>}{\sqrt{2}}
  	  \] \pause
  
  We apply the Hadamard gate to obtain
  	\[
  			|x_{k+1}\>
  	\] \pause
  	
  The controlled phase shifts enable us to reduce the problem of determining each bit to distinguishing between $|+\>$ and $|-\>$ (deterministic Hadamard test)

\end{frame}

%%%%%%%%%%%%%%%%%%%%%%%%%%%%%%%%%%%%%%%%%%%%%%%%%%%%%%%%%%%%%%%%%%%%%%%%%%%%

%% Special case

%%%%%%%%%%%%%%%%%%%%%%%%%%%%%%%%%%%%%%%%%%%%%%%%%%%%%%%%%%%%%%%%%%%%%%%%%%%%

\begin{frame}{Special case: exact $n$-bit binary fraction}

Assume that $\varphi$ is an exact $n$-bit binary fraction, i.e., $\varphi=0.x_1 \ldots x_{n-1} x_n$

\bigskip
\hspace{1cm}
\Qcircuit @C=1.0em @R=0.7em {
\lstick{|0\>}& \qw & \gate{H}                 & \qw      & \qw      & \qw &        &        & \qw & \ctrl{4} & \qw & \multigate{3}{F^\dagger} & \qw & \rstick{|x_n\>}\\
\lstick{|0\>}& \qw & \gate{H}                 & \qw      & \ctrl{3} & \qw &        &        & \qw & \qw      & \qw & \ghost{F^\dagger}        & \qw & \rstick{|x_{n-1}\>}\\
             &     & \raisebox{0.5em}{\vdots} &          &          &     & \cdots &        &     &          &     &         &     & \\
\lstick{|0\>}& \qw & \gate{H}                 & \ctrl{1} & \qw      & \qw &        &        & \qw & \qw      & \qw & \ghost{F^\dagger}        & \qw & \rstick{|x_1\>}\\
%
\lstick{|\psi\>} & \qw & {/} \qw  & \gate{U^{2^0}} & \gate{U^{2^1}} & \qw & & & \qw & \gate{U^{2^{n-1}}} & \qw & \qw & \qw & \rstick{|\psi\>}
}

\bigskip
$\Rightarrow$ The measurment of the qubits yields the bits $x_n, x_{n-1},\ldots, x_1$ deterministically

\end{frame}

%%%%%%%%%%%%%%%%%%%%%%%%%%%%%%%%%%%%%%%%%%%%%%%%%%%%%%%%%%%%%%%%%%%%%%%%%%%%%%%%%%%%%%%%%%%%%%%%%%%%%%%%%%%

%% General case

%%%%%%%%%%%%%%%%%%%%%%%%%%%%%%%%%%%%%%%%%%%%%%%%%%%%%%%%%%%%%%%%%%%%%%%%%%%%%%%%%%%%%%%%%%%%%%%%%%%%%%%%%%%

\begin{frame}{General case: arbitrary eigenphases}

Let $\varphi$ be arbitrary

\bigskip
Unless $\varphi$ is an exact $n$-bit fraction, the application of the inverse quantum Fourier transform
\bigskip
\[
		F^\dag |\varphi\> 
\]
produces a {\color{blue}superposition} of $n$-bit strings
 
\end{frame}

%%%%%%%%%%%%%%%%%%%%%%%%%%%%%%%%%%%%%%%%%%%%%%%%%%%%%%%%%%%%%%%%%%%%%%%%%%%%%%%

%% Geometric summation

%%%%%%%%%%%%%%%%%%%%%%%%%%%%%%%%%%%%%%%%%%%%%%%%%%%%%%%%%%%%%%%%%%%%%%%%%%%%%%%

\begin{frame}{Geometric summation}

\begin{lemma}
We have
\begin{eqnarray*}
\sum_{y=0}^{N-1} e^{2\pi i \theta y} & = & N \mbox{~for $\theta=0$} \\ \pause
\sum_{y=0}^{N-1} e^{2\pi i \theta y} & = & \frac{1-e^{2\pi i N \theta}}{1-e^{2\pi i \theta}} \mbox{~for $\theta\in (0,1)$}
\end{eqnarray*}
\end{lemma}

\bigskip
Assume that $\theta=\frac{x}{N}$ for some $x\in[0,N-1]$

\pause

\bigskip
$\Rightarrow$ We have

\[
\sum_{y=0}^{N-1} e^{2\pi i \frac{x}{N} y} = N \, \delta_{x,0}
\]

\bigskip

\end{frame}

%%%%%%%%%%%%%%%%%%%%%%%%%%%%%%%%%%%%%%%%%%%%%%%%%%%%%%%%%%%%%%%%%%%%%%%%%%%

%% Probability of measuring a bit string

%%%%%%%%%%%%%%%%%%%%%%%%%%%%%%%%%%%%%%%%%%%%%%%%%%%%%%%%%%%%%%%%%%%%%%%%%%%

\begin{frame}{Probability of obtaining a certain estimate}

\begin{lemma}
Let $x=\sum_{k=1}^n x_i 2^{n-i}$ and $\varphi_x := 0.x_1 x_2 \ldots x_n = \frac{x}{2^n}$ be the corresponding $n$-bit fraction \pause

\bigskip
The probability of obtaining the estimate $\varphi_x$ is 
\bigskip
\begin{eqnarray*}
	\Pr(x) 
	& = &
	\frac{1}{2^{2n}}
	\frac{\sin^2 \big(2^n \, \pi \, (\varphi - \varphi_x)\big)}{\sin^2\big( \pi\, (\varphi - \varphi_x)\big)}
\end{eqnarray*}
\end{lemma}
\end{frame}

%%%%%%%%%%%%%%%%%%%%%%%%%%%%%%%%%%%%%%%%%%%%%%%%%%%%%%%%%%%%%%%%%%%%%%%%%%%%%%%

%% Probability distribution

%%%%%%%%%%%%%%%%%%%%%%%%%%%%%%%%%%%%%%%%%%%%%%%%%%%%%%%%%%%%%%%%%%%%%%%%%%%%%%%

\begin{frame}{Examples of probability distributions for different $\varphi$}

%\frame{
%\movie[width=4in,height=3in,poster]{}{phase.avi}

\begin{center}
\scalebox{.7}{
\multiinclude[<+-| handout:+->][format=png]{PE}
}
\end{center}

\end{frame}

%%%%%%%%%%%%%%%%%%%%%%%%%%%%%%%%%%%%%%%%%%%%%%%%%%%%%%%%%%%%%%%%%%%%%%%%%%%%%%%

%% Proof

%%%%%%%%%%%%%%%%%%%%%%%%%%%%%%%%%%%%%%%%%%%%%%%%%%%%%%%%%%%%%%%%%%%%%%%%%%%%%%%

\begin{frame}{Probability of obtaining a certain estimate}

\begin{proof}
The probability of obtaining the estimate $\varphi_x$ is 
\begin{eqnarray*}
&   & 
\Pr(x) \\
& = & 
|\< x| F^\dag |\varphi\> |^2 \\ \pause
& = & 
| \< \varphi_x| \varphi\> |^2 \\ \pause
& = &
\frac{1}{2^{2n}} 
\big|
\sum_{y = 0}^{2^n-1} e^{2 \pi i \, (\varphi -
\varphi_x) \, y }
\big|^2 \pause \quad \mbox{geometric summation} \\ \pause
& = &
\frac{1}{2^{2n}}
\Big|
\frac{1 - e^{2 \pi i (2^n (\varphi - \varphi_x))}}{1 -
e^{2 \pi i (\varphi - \varphi_x)}}
\Big|^2 \pause \quad |1 - e^{i 2\theta}| =| e^{-i \theta} -
e^{i\theta}| = 2|\sin\theta| \\ \pause
& = &
\frac{1}{2^{2n}}
\frac{\sin^2 (2^n\, \pi (\varphi - \varphi_x))}{\sin^2(\pi(\varphi - \varphi_x))}
\end{eqnarray*}
\end{proof}

\end{frame}

%%%%%%%%%%%%%%%%%%%%%%%%%%%%%%%%%%%%%%%%%%%%%%%%%%%%%%%%%%%%%%%%%%%%%%%%%%%%%%%%%%

%% Success probability

%%%%%%%%%%%%%%%%%%%%%%%%%%%%%%%%%%%%%%%%%%%%%%%%%%%%%%%%%%%%%%%%%%%%%%%%%%%%%%%%%%

\begin{frame}{Lower bound on success probability}

\begin{theorem}
Let $x$ be such that $\frac{x}{2^n} \le \varphi < \frac{x+1}{2^n}$

\bigskip
The probability of returning one of the two closest $n$-bit fractions $\varphi_x$ and $\varphi_{x+1}$ is at least $\frac{8}{\pi^2}$
\end{theorem}

\end{frame}

%%%%%%%%%%%%%%%%%%%%%%%%%%%%%%%%%%%%%%%%%%%%%%%%%%%%%%%%%%%%%%%%%%%%%%%%%%%%%%%%%%%%%

%% Proof continued

%%%%%%%%%%%%%%%%%%%%%%%%%%%%%%%%%%%%%%%%%%%%%%%%%%%%%%%%%%%%%%%%%%%%%%%%%%%%%%%%%%%%%

\begin{frame}{Proof of lower bound on success probability}

\begin{eqnarray*}
\Pr(\mbox{success}) & := &
\Pr(x) + \Pr(x+1) \\
& = &
\frac{1}{2^{2n}}
\left( \big| 
\sum_{y = 0}^{{2^n}-1} 
e^{2 \pi i (\varphi - \varphi_x)y}
\big|^2 
+ 
\big|\sum_{y = 0}^{{2^n}-1}
e^{2 \pi i (\varphi - \varphi_x)y}
\big|^2\right)
\end{eqnarray*}

\pause

\medskip
This function attains its minimum at $\varphi = \frac{1}{2} (\varphi_x + \varphi_{x+1})$ \pause \quad $\Rightarrow$
\begin{eqnarray*}
\Pr(\mbox{success}) 
& \geq &
\frac{2}{2^{2n}}
\left(
\big|
\sum_{y = 0}^{{2^n}-1} e^{2 \pi i \frac{y}{2^{n+1}}}
\big|^2
\right) \\ \pause
& \geq &
\frac{2}{2^{2n}}
\frac{4}{4 \sin^2(\frac{\pi}{2^{n+1}})} \\ \pause
& \geq &
\frac{8}{\pi^2}
\end{eqnarray*}

The last inequality follows from $\frac{1}{|\sin{\theta}|^2} \geq \frac{1}{|\theta|^2}$
\end{frame}

%%%%%%%%%%%%%%%%%%%%%%%%%%%%%%%%%%%%%%%%%%%%%%%%%%%%%%%%%%%%%%%%%%%%%%%%%%%%%%%%

%% Summary of phase estimation

%%%%%%%%%%%%%%%%%%%%%%%%%%%%%%%%%%%%%%%%%%%%%%%%%%%%%%%%%%%%%%%%%%%%%%%%%%%%%%%%

\begin{frame}{Summary of phase estimation}

\bigskip
We are given a unitary $U$ and an eigenvector $|\psi\>$ of $U$ with unknown eigenphase $\varphi$

\bigskip
We obtain an estimate $\hat{\varphi}$ such that
\[
   \Pr\big(|\hat{\varphi} - \varphi| \le \frac{1}{2^n}\big) \ge \frac{8}{\pi^2}
\]

\medskip
To do this, we need invoke each of the controlled $U$, $U^2$,\ldots,$U^{2^{n-1}}$ gates once \pause

\bigskip
We can boost the success probability to $1-\epsilon$ by repeating the above algorithm $O(\log(1/\epsilon))$ times and outputting the median of the outcomes

\end{frame}

%%%%%%%%%%%%%%%%%%%%%%%%%%%%%%%%%%%%%%%%%%%%%%%%%%%%%%%%%%%%%%%%%%%%%%%%%%%%%%%%

%% Phase estimation

%%%%%%%%%%%%%%%%%%%%%%%%%%%%%%%%%%%%%%%%%%%%%%%%%%%%%%%%%%%%%%%%%%%%%%%%%%%%%%%%

\begin{frame}{Phase estimation applied to superpositions of eigenstates}

\bigskip
We are given a unitary $U$ with eigenvectors $|\psi_i\>$ and corresponding eigenphases $\varphi_i$

\bigskip
Let 
\[
|\psi\> = \sum_i \alpha_i |\psi_i\>
\]

\pause

\bigskip
What happens if we apply phase estimation to $|0\>^{\otimes n} \otimes |\psi\>$?

\pause

\bigskip
After the $n$ phase kick-backs due to $U^{2^0}$, $U^{2^1}$, \ldots $U^{2^{n-1}}$, we obtain

\[
\sum_i \alpha_i |\varphi_i\> \otimes |\psi_i\>
\]

\pause

\bigskip
After applying the inverse quantum Fourier transform, we obtain

\[
\sum_i \alpha_i |\tilde{x}_i\> \otimes |\psi_i\>
\]

where $|\tilde{x}_i\>$ denotes a superpositions of $n$-bit estimates of $\varphi_i$

\end{frame}
